\documentclass[fullpage,twoside,english,12pt]{book}

\begin{document}

\section*{Abstract}
In this talk I would like to present you the main results of my PhD
thesis. They encompass the formal definition of mutual simulation of
Hamiltonians and the development of a theory investigating the
computational resources required for simulating Hamiltonians.

Mutual simulation of Hamiltonians builds upon a control theoretic
model of quantum computers. In this model unitary transformations are
implemented by performing an external control operation, letting the
system evolve for a certain time, and repeating both steps several
times. We study the ability of the system to implement transformations
representing the time evolution of another quantum system. Considering
the cost of simulating the time evolution of a quantum system by a
given one defines in a certain sense the comparison of the
``computational power'' of different systems. The determination of the
cost in general is a computationally hard problem. But it can be
solved to a large extent if we consider infinitesimal time intervals
only. This leads to our definition of mutual simulation of
Hamiltonians.

We derive lower bounds on the simulation complexity. An important tool
for obtaining these bounds is the characterization of Hamiltonians by
graphs; the graphs represent the coupling topology of the
Hamiltonians, that is, which pairs of subsystems interact. Our bounds
show that the simulation complexity depends strongly on the properties
of these graphs. We develop optimal simulation algorithms based on
methods of graph theory, representation theory of finite groups,
design theory, and convex optimization theory. Furthermore, we show
how the elementary quantum gates of the quantum circuit model can be
implemented within the control theoretic model. This shows an
interesting connection between the quantum circuit model and the
control theoretic model.

We show as an interesting application of our simulation methods that
they allow an efficient realization of adiabatic quantum computing in
the control theoretic model. It is important that only physically
realistic Hamiltonians (that is, with a local coupling topology) are
used as computational resources. We construct such Hamiltonians whose
energy minima encode the solutions of the NP-complete problem ``max
independent set''. Due to their special structure these Hamiltonians
can be simulated efficiently by Ising-interactions on a
two-dimensional lattice.



\end{document}
